\chapter{1999年上海卷（文）}
\section{填空题}

\begin{problem}
\(\operatorname{tg}\left[\arccos\frac{\sqrt{2}}{2}-\frac{\pi}{6}\right]= \) \fillinblank{}.

\end{problem}

\begin{problem}
函数 \(f(x)=\log_{2}x+1(x\geqslant4) \) 的反函数 \(f^{-1}(x) \) 的定义域是 \fillinblank{}.

\end{problem}

\begin{problem}
在 \((x^{3}+\frac{2}{x^{2}})^{5} \) 的展开式中，含 \(x^{5} \) 项的系数为 \fillinblank{}.

\end{problem}

\begin{problem}
在\(\triangle ABC\)中，若\(\angle B=30^{\circ}, AB=2\sqrt{3}, AC=2\)，则\(\triangle ABC\)的面积 \(S\) 是 \fillinblank{}.

\end{problem}

\begin{problem}
若平移坐标系，将曲线方程 \(y^{2} + 4x - 4y - 4 = 0 \) 化为标准方程，则坐标原点应移到点 \(O'(\fillinblank{}, \fillinblank{}) \) .

\end{problem}

\begin{problem}
函数 \(y = 2\sin x \cos x - 2\sin^{2}x + 1 \) 的最小正周期 \(T=\) \fillinblank{}.

\end{problem}

\begin{problem}
函数 \(y = 2\sin(2x + \frac{\pi}{6})(x \in [-\pi, 0]) \) 的单调递减区间是 \fillinblank{}.

\end{problem}

\begin{problem}
某工程的工序流程图如下(工时单位: 天),

现已知工程总时数为 10 天，则工序 \(c\) 所需工时为

天 \fillinblank{}.

\end{problem}

\begin{problem}
\(\frac{\log_{3}2}{\log_{27}64}= \) \fillinblank{}.

\end{problem}

\begin{problem}
在等差数列 \(\{a_{n}\} \) 中满足 \(3a_{4}=7a_{7} \) ，且 \(a_{1}>0, S_{n} \) 是数列 \(\{a_{n}\} \) 前 \(n\) 项的和. 若 \(S_{n} \) 取得最大值，则 \(n=\) \fillinblank{}.

\end{problem}

\begin{problem}
若以连续掷两次骰子分别得到的点数 \(m\)、\(n\) 作为点 \(P\) 的坐标，则点 \(P\) 落在圆 \(x^{2} + y^{2} = 16 \) 内的概率是 \fillinblank{}.

\end{problem}

\begin{problem}
若四面体各棱的长是 1 或 2，且该四面体不是正四面体，则其体积的值是 (只需写出一个可能的值) \fillinblank{}.

\end{problem}

\section{单选题}

\begin{problem}
直线 \(y=\frac{\sqrt{3}}{3}x \) 绕原点按逆时针方向旋转 \(30^{\circ} \) 后所得直线与圆 \((x-2)^{2}+y^{2}=3 \) 的位置关系是（\quad）

\choices
    {直线过圆心}
    {直线与圆相交，但不过圆心}
    {直线与圆相切}
    {直线与圆没有公共点}

\end{problem}

\begin{answer}
B
\end{answer}

\begin{problem}
下列以 \(t\) 为参数的参数方程所表示的曲线中，与方程 \(xy = 1\) 所表示的曲线完全一致的是（\quad）

\choices
    {\(\left\{\begin{aligned}x&=\frac{1}
    {2}\\ y&=t-\frac{1}
    {2}\end{aligned}\right. \)}
    {\(\left\{\begin{aligned}x&=|t|\\ y&=\frac{1}{|t|}\end{aligned}\right. \)}
 {\(\left\{\begin{aligned}x&=\cos t\\ y&=\sec t\end{aligned}\right. \)}
 {\(\left\{\begin{aligned}x&=\operatorname{tg}t\\ y&=\operatorname{ctg}t\end{aligned}\right. \)}

\end{problem}

\begin{answer}
D
\end{answer}

\begin{problem}
若 \(a < b < 0\)，则下列结论中正确的是（\quad）

\choices
    {不等式 \(\frac{1}
    {a} > \frac{1}
    {b}\) 和 \(\frac{1}
    {|a|} > \frac{1}{|b|}\) 均不能成立}
 {不等式 \(\frac{1}{a - b} > \frac{1}{a}\) 和 \(\frac{1}{|a|} > \frac{1}{|b|}\) 均不能成立}
 {不等式 \(\frac{1}{a - b} > \frac{1}{a}\) 和 \((a + \frac{1}{b})^2 > (b + \frac{1}{a})^2\) 均不能成立}
 {不等式 \(\frac{1}{|a|} > \frac{1}{|b|}\) 和 \((a + \frac{1}{b})^2 > (b + \frac{1}{a})^2\) 均不能成立}

\end{problem}

\begin{answer}
A
\end{answer}

\begin{problem}
设有直线 \(m\)、\(n\) 和平面 \(\alpha \) 、 \(\beta \) ，则在下列命题中，正确的是（\quad）

\choices
    {若 \(m // n, m \subset \alpha, n \subset \beta\)，则 \(\alpha // \beta\)}
    {若 \(m \perp \alpha, m \perp n, n \subset \beta\)，则 \(\alpha // \beta\)}
    {若 \(m // n, n \perp \beta, m \subset \alpha\)，则 \(\alpha \perp \beta\)}
    {若 \(m // n, m \perp \alpha, n \perp \beta\)，则 \(\alpha \perp \beta\)}

\end{problem}

\begin{answer}
A
\end{answer}

\section{解答题}

\begin{problem}
设集合 \(A = \{x\mid |x - a| < 2\}, B = \{x\mid \frac{2x - 1}{x + 2} < 1\}\). 若 \(A \subseteq B\)，求实数 \(a\) 的取值范围.

\end{problem}

\begin{problem}
设正数数列 \(\{a_{n}\} \) 为一等比数列，且 \(a_{2}=4, a_{4}=16 \) ，求 \(\displaystyle\lim_{n\to\infty}\frac{\lg a_{n+1}+\lg a_{n+2}+\cdots+\lg a_{2n}}{n^{2}} \) .

\end{problem}

\begin{problem}
设复数 \(z\) 满足 \(4z + 2\overline{z} = 3 + \sqrt{3} + \i \) ，\(w = \sin\theta - \i\cos\theta (\theta \in R) \) ，求 \(z\) 的值和 \(|z - w| \) 的取值范围.

\end{problem}

\begin{problem}
如图，在四棱锥 \(P\)—\(ABCD\) 中，底面 \(ABCD\) 是一直角梯形， \(\angle BAD = 90^{\circ} \) ，\(AD // BC\)，\(AB = BC = a\)，\(AD = 2a\)，且 \(PA \perp \) 底面 \(ABCD\)，\(PD\) 与底面成 \(30^{\circ} \) 角. （1）若 \(AE \perp PD \) ，\(E\) 为垂足，求证： \(BE \perp PD \) ;

（2）求异面直线 \(AE\) 与 \(CD\) 所成角的大小(用反三角函数表示).

\end{problem}

\begin{problem}
平地上有一条水沟，沟沿是两条长 100 米的平行线段，沟宽 \(AB\) 为 2 米，与沟沿垂直的平面与沟的交线是一段抛物线，抛物线顶点为 \(O\)，对称轴与地面垂直，沟深 1.5 米，沟中水深 1 米. （1）求水面宽； （2）如图所示形状的几何体积为柱体.已知柱体的体积为底面积乘以高.问沟中的水有多少立方米？ （3）若要把这条水

沟改挖(不准填土)成截面为等腰梯形的沟，使沟的底面与地面平行，则改挖后的沟底宽为多少米时，所挖的土最少\(\perp\)

\end{problem}

\begin{problem}
设椭圆 \(l\) 的方程为 \(\frac{(x+c)^{2}}{b^{2}}+\frac{(y+c)^{2}}{a^{2}}=1 \) ，其中 \(a>b>0\)，\(c\) 为椭圆的半焦距. （1）将椭圆 \(l\) 绕它的上焦点按逆时针方向旋转 \(\frac{\pi}{2} \) ，求所得椭圆 \(L\) 的中心的坐标、长轴长、短轴长及其方程； （2）若椭圆 \(L\) 与曲线 \(y=\frac{1}{x} \) 在第一象限内只有一个公共点 \(P\)，试用 \(a\) 表示点 \(P\) 的坐标； （3）设 \(A\)、\(B\) 是椭圆 \(L\) 的两个焦点，当 \(a\) 变化时，求 \(\triangle ABP \) 的面积函数 \(S(a) \) 的值域.
\end{problem}

